\subsection{二項係数で多項式を表す}\label{節：二項係数で多項式を表す}
平方和を,今度は\sectionref*{節：二項定理}の二項係数で求める.
ここでは$n,k,N$は非負整数とし,
\[
  \binom nk=\frac{n!}{k!(n-k)!}\quad(0\le k\le n)
\]
とする.$0!=1$, $k>n$では$\binom nk=0$と約束する.

\bookstep{例題：二項係数で平方和を求める}
$n^2=n(n-1)+n$だから,
\[
  n^2=2\binom n2+\binom n1
\]
と表せる.
部品の差は
\[
  \binom{n+1}{3}-\binom n3=\binom n2\quad(n\ge0)
\]
となる.$n\ge3$では階乗表示を展開して確かめられ,$n=0,1,2$も上の規約で成り立つ.
同じように,$\binom{n+1}{2}-\binom n2=\binom n1$である.
$j=0$から$N-1$まで足すと途中が消え,従って
\BookMathLayout{\[
    \sum_{j=0}^{N-1}j^2=2\binom N3+\binom N2
    =\frac{N(N-1)(2N-1)}6.
  \]}{\[
    \begin{aligned}
      \sum_{j=0}^{N-1}j^2&=2\binom N3+\binom N2\\
      &=\frac{N(N-1)(2N-1)}6.
    \end{aligned}
  \]}
$N=n+1$とすれば,先ほどの$1^2+\cdots+n^2$の公式に戻る.
二項係数なら,添字を一つ下げるだけで差を表せる.

\bookstep{一般の添字で確かめる}
$n+1$個から$k$個を選ぶとき,特定の一個を選ぶかどうかで分ければ,
\[
  \binom{n+1}{k}=\binom nk+\binom n{k-1}\quad(k\ge1)
\]
である.従って
\[
  \Delta\binom nk=\binom n{k-1}\quad(k\ge1)
\]
となる.$\Delta$は前進の差で,$\binom n0=1$の差は別に零となる.

逆に差を足すと,途中の項が消えるので,
\[
  \sum_{j=0}^{N-1}\binom jk=\binom N{k+1}
  \quad(N,k\ge0)
\]
となる.$N=0$では空和$0$,$k=0$では和が$N$となり,右辺と一致する.
\BookFigBinomialDifference

\begin{bookpoint}
  \textbf{式の部品を選ぶ}\par
  目的の操作が簡単になる\bookterm{basis}{基底}を選ぶ.
  \sectionref*{節：隣接3項型}の見方が,多項式の表示にも現れた.
\end{bookpoint}
実際,$\binom n0,\ldots,\binom nd$はそれぞれ次数が$0,\ldots,d$で,
最高次係数が非零である.高い次数から係数を合わせれば,
$d$次以下の多項式を一意にその\bookterm{combination}{線形結合}で表せる.
三乗和は\bookproblemref{practice-connections}{7}{接続演習問7}で試そう.
